\documentclass[11pt]{article}
\usepackage[utf8]{inputenc}
\usepackage{amsmath, amssymb}
\usepackage{geometry}
\usepackage{listings}
\usepackage{xcolor}
\usepackage{tabularx}
\usepackage{booktabs}
\usepackage{hyperref}

\geometry{a4paper, margin=1in}

\definecolor{codegray}{rgb}{0.5,0.5,0.5}
\definecolor{codepurple}{rgb}{0.58,0,0.82}
\definecolor{backcolour}{rgb}{0.95,0.95,0.92}

\lstset{
    language=Python,
    backgroundcolor=\color{backcolour},
    commentstyle=\color{codegray},
    keywordstyle=\color{magenta},
    stringstyle=\color{codepurple},
    basicstyle=\ttfamily\small,
    breaklines=true,
    showstringspaces=false
}

\title{Game Theory Lab: Evolutionary Strategies \& Spatial Chaos \\ \large Problem Set \& Implementation Guide}
\author{VB}
\date{\today}

\begin{document}

\maketitle

\section*{Problem 0: Start the engine}

Consider the files given.
\begin{center}
\renewcommand{\arraystretch}{1.5}
\begin{tabularx}{\textwidth}{@{} >{\bfseries\raggedright\arraybackslash}p{5.5cm} X @{}}
\toprule
\textcolor{blue}{Filename} & \textcolor{blue}{Purpose} \\ \midrule
\multicolumn{2}{@{}l}{\textbf{Direct Interaction}} \\ \midrule
\lstinline|implementation_tasks.py| & \textcolor{red}{Here you write your code.} Contains the core evolutionary algorithms, fitness evaluations, and your custom Strategy classes. \\
\lstinline|lab_dashboard.py| & \textcolor{red}{This file you run.} A dedicated educational sandbox for exploring 1v1 matches, Axelrod tournaments, and Nowak \& May spatial grids. \\
\lstinline|levels/level_*.py| & Graphical representation for each specific problem. \textbf{Can be run standalone.} \\ \midrule
\multicolumn{2}{@{}l}{\textbf{Under the Hood}} \\ \midrule
\lstinline|levels/base_level.py| & Base class for graphical debuggers, managing Tkinter layouts and Matplotlib interaction.\\
\bottomrule
\end{tabularx}
\end{center}

Run the file \lstinline|lab_dashboard.py|. You will see a window with a menu of evolutionary game theory levels.

Now use any editor to open the file \lstinline|implementation_tasks.py|. You should see the API for functions you are expected to implement. You can add your own \lstinline|unittest.TestCase| classes at the bottom and run \lstinline|python implementation_tasks.py| to execute them.

\subsubsection*{Self-Testing (Optional)}
At the bottom of \lstinline|implementation_tasks.py|, you will find a block that looks like this:
\begin{lstlisting}[language=Python]
if __name__ == '__main__':
    import unittest
    # Add your own unittest.TestCase classes here...
    unittest.main(verbosity=2)
\end{lstlisting}
This allows you to write your own automated unit tests to verify your code before running the visual dashboard. If you'd like to test your functions with specific inputs, simply define a class inheriting from \lstinline|unittest.TestCase| above the \lstinline|unittest.main()| call. For example, if you wanted to test simple math:
\begin{lstlisting}[language=Python]
class TestMyMath(unittest.TestCase):
    def test_addition(self):
        result = 2 + 2
        self.assertEqual(result, 4, "Addition is broken!")
\end{lstlisting}
You can run these tests anytime by executing \lstinline|python implementation_tasks.py| in your terminal. This is completely optional, but highly recommended for debugging tricky edge cases.
\section*{Problem 1: Custom Strategies \& 1v1 Matches}

\subsubsection*{Mathematical part}
\textbf{Given:} A standard $2 \times 2$ symmetric game, like the Prisoner's Dilemma, defined by a payout matrix with entries for Temptation ($T$), Reward ($R$), Punishment ($P$), and Sucker ($S$). \\
\textbf{Derive:} The total expected payout for two strategies playing against each other for $N$ rounds. Strategies may use the entire history of the game to make decisions (e.g. \textit{Tit-for-Tat} or \textit{Grim Trigger}).

\subsubsection*{Implementation part}
\textbf{Implement:}
\begin{itemize}
    \item \textbf{Fundamental Strategies:} You are required to implement four foundational strategies: \textit{AlwaysCooperate}, \textit{AlwaysDefect}, \textit{TitForTat}, and \textit{GrimTrigger}. You will find the class stubs inside \lstinline|implementation_tasks.py|. You must complete their \lstinline|get_action(self, my_history, opponent_history)| methods to return \lstinline|'C'| or \lstinline|'D'|.
    \item \textbf{Custom Strategies:} After implementing the basics, you are highly encouraged to design your own custom evolutionary strategies by inheriting from \lstinline|BaseStrategy|.
    \item \textbf{Match Engine:} Implement \lstinline|play_match(strat_A, strat_B, matrix, rounds)|. This function should reset both strategies, iterate $N$ rounds, ask each strategy for its action, record the history, and tally the total payout from the \lstinline|matrix|.
\end{itemize}

\section*{Problem 2: Spatial Fitness Evaluation}

\subsubsection*{Mathematical part}
\textbf{Given:} A 2D spatial lattice of agents. \\
\textbf{Derive:} The fitness of an agent at a given position. In our model, an agent's fitness is the sum of its payouts playing a 1-shot game against all 8 of its neighbors (a Moore neighborhood).

\subsubsection*{Implementation part}
\textbf{Implement:}
\begin{itemize}
    \item \lstinline|compute_neighborhood_fitness(grid, r, c, matrix, rounds)|: Iterate through the $3 \times 3$ subgrid surrounding the target cell $(r, c)$. Make sure to apply toroidal (wrap-around) boundary conditions using modulo arithmetic. Exclude the agent itself from the calculation. Return the total accumulated payout.
\end{itemize}

\section*{Problem 3: Spatial Chaos (Imitate-the-Best)}

\subsubsection*{Mathematical part}
\textbf{Given:} A 2D spatial lattice of agents and their fitnesses. \\
\textbf{Derive:} The state of the lattice in the next generation. We use a simple deterministic rule: every cell looks at its 8 neighbors and itself (9 cells total). In the next generation, it adopts the strategy of whoever had the highest fitness in that local neighborhood. \\

\textit{Note on Parameters:} As demonstrated by Nowak and May (Nature, 1992), this simple deterministic rule produces chaotic, beautiful, fractal-like patterns only under very specific parameter regimes. If you set $R=1, P=0, S=0$, the dynamics are entirely controlled by the Temptation parameter $T$ (often denoted $b$):
\begin{itemize}
    \item \textbf{$1.8 < T < 2.0$ (The Chaotic Window):} Defectors are strong enough to invade, but Cooperators can band together to survive and reclaim territory. This generates infinite, shifting ''kaleidoscopes'' and ''gliders''. Try $T=1.9$ with a single Defector in the center!
    \item \textbf{$T < 1.8$:} Defectors are too weak and will generally die out.
    \item \textbf{$T > 2.0$:} Defectors are too strong and will overrun the entire grid, leaving a dead world of 0 payout.
\end{itemize}

\subsubsection*{Implementation part}
\textbf{Implement:}
\begin{itemize}
    \item \lstinline|update_grid_deterministic(grid, matrix, rounds)|: Precompute the fitness of every agent on the grid using \lstinline|compute_neighborhood_fitness|. Then, for each cell, find the local maximum fitness among its 8 neighbors and itself. Return a \textbf{new} grid populated with newly instantiated strategy objects corresponding to the local winners. Do not modify the grid in-place, as this causes tearing artifacts!
\end{itemize}

\end{document}
